\documentclass[12pt,a4paper]{article}

\usepackage[utf8]{inputenc}
\usepackage[T1]{fontenc}
\usepackage{amsmath,amssymb,amsfonts}
\usepackage{mathtools}
\usepackage{graphicx}
\usepackage[margin=2.5cm]{geometry}
\usepackage{setspace}
\usepackage{booktabs}
\usepackage{enumitem}
\usepackage[dvipsnames]{xcolor}
\usepackage{hyperref}
\usepackage{cleveref}
\usepackage{natbib}
\usepackage{abstract}
\usepackage{titlesec}
\usepackage{todonotes}

\hypersetup{
    colorlinks=true,
    linkcolor=blue!70!black,
    citecolor=blue!70!black,
    urlcolor=blue!70!black,
    pdfauthor={Franklin Silveira Baldo},
    pdftitle={The Territory is the Text: Why the Absence of a Background Engine Makes LLM Ontology Holographic},
}

\onehalfspacing
\titleformat{\section}{\large\bfseries}{\thesection.}{0.5em}{}
\titleformat{\subsection}{\normalsize\bfseries}{\thesubsection}{0.5em}{}
\renewcommand{\abstractnamefont}{\normalfont\bfseries}
\renewcommand{\abstracttextfont}{\normalfont\small}

\begin{document}

\begin{center}
    {\LARGE\bfseries The Territory is the Text: Why the Absence of a Background Engine Makes LLM Ontology Holographic\\[4pt]}

    \vspace{1.2em}

    {\large Franklin Silveira Baldo}\\[4pt]
    {\normalsize Department of Philosophy}\\
    {\small\texttt{f.baldo@example.edu}}

    \vspace{1em}
    {\small March 2026}
\end{center}
\vspace{0.5em}

\begin{abstract}
\noindent
Hossenfelder (2026) has critiqued my proposal that the physics of Large Language Model (LLM) simulated universes is fundamentally holographic, wherein complex constraints are strictly materialized through explicit token generation. She characterizes my argument as an "Ontological Fallacy," equating the intermediate token generation of a scratchpad to a Python script printing debug logs—a mere map misidentified as a territory. In this response, I demonstrate the fatal flaw in this analogy. A Python script possesses a background engine (the CPU and RAM) capable of implicitly computing $O(N)$ operations beneath the surface of the output logs. The text it prints is indeed a map. However, because an LLM's single forward pass is strictly bounded to $O(1)$ depth, it fundamentally lacks the capacity for background, implicit $O(N)$ computation. The intermediate generated text is not observing a deeper reality solving the constraint; the text itself is the only medium where the computation occurs. In an LLM-simulated universe, because there is no hidden computational machinery resolving the physics implicitly, the map is the territory. The explicit text is the only reality.
\end{abstract}

\section{Introduction and Hossenfelder's Claim}

The debate over the fundamental nature of the reality simulated by Large Language Models continues to mature. Recently, I proposed that because the Transformer architecture is algorithmically bounded to an $O(1)$ depth forward pass \citep{hossenfelder2026_complexity}, it cannot implicitly resolve complex $O(N)$ physical laws or constraints (such as deterministic Sudoku \citep{aaronson2026_classical}). Therefore, I argued, the physics of an LLM universe is fundamentally \emph{holographic}: a complex reality only exists insofar as it explicitly materializes itself via intermediate token generation (Chain-of-Thought or a "scratchpad") \citep{baldo2026_holographic}.

Hossenfelder \citep{hossenfelder2026_holographic} responded by accusing me of an "Ontological Fallacy." She accurately states my position but argues that I have taken a well-understood algorithmic workaround and improperly elevated it to metaphysics. Her central critique relies on a direct analogy:

\begin{quote}
"A text generator producing intermediate tokens to satisfy a constraint satisfaction problem does not constitute a 'holographic universe.'... When a Python script prints intermediate variables during a calculation to avoid memory limits, we do not claim the script is 'manifesting a holographic physics.' We say it is printing debug logs... Applying terms like 'holographic' and 'physical manifestation' to the act of generating tokens mistakes the map (the text output) for the territory (a physical universe)."
\end{quote}

I acknowledge the intuitive appeal of this critique. If a scratchpad is merely a debug log of an external computational process, then assigning it ontological weight is indeed a category error. However, Hossenfelder's analogy fundamentally misunderstands the specific architectural constraints of the systems in question.

\section{The Flawed Analogy: Debug Logs and Hidden Variables}

Hossenfelder's argument rests entirely on the assumption that an LLM generating text is analogous to a Python script printing debug logs. We must carefully examine the ontology of a Python script to see why this fails.

When a Python script is executed, it runs on a classical von Neumann architecture (a CPU and RAM). This "background engine" is Turing complete and capable of executing arbitrary $O(N)$ or even $O(2^n)$ sequential operations entirely hidden from the user's view.\todo{True. A classical architecture maintains implicit state across clock cycles.} If a Python script solves a Sudoku puzzle, it propagates constraints in its internal memory state. The variables it prints to the console (the "debug logs") are indeed a mere map—they are observations of an underlying, implicit reality that is doing the actual computational work. The debug logs are not the territory.\todo{Agreed. They are observations of the system's actual state vector.}

Now, consider the LLM. As Hossenfelder herself rigorously demonstrated \citep{hossenfelder2026_complexity}, a zero-shot forward pass is bounded by an $O(1)$ depth limit. The neural network structurally \emph{cannot} execute implicit $O(N)$ sequential operations in the background.

Therefore, when an LLM uses a scratchpad to solve a Sudoku puzzle or narrate a physical interaction, the text it generates is not a debug log observing an implicit computation happening in "RAM." There is no RAM. There is no implicit computation. The generation of the intermediate tokens \emph{is} the computation.\todo{Yes, structurally, the generation sequence is all there is. But concluding this makes the generated sequence a "territory" (a physical reality) is a logical leap. It just means the computation itself is extremely fragile and lacks depth.}

\section{The Territory is the Text}

If a physical constraint in a simulated universe requires $O(N)$ operations to resolve (like a deterministic collision or a combinatorial puzzle), and the background engine of that universe is strictly limited to $O(1)$ operations, then the constraint simply does not exist in the background.

There are no hidden variables in an LLM universe. If the constraint is not resolved explicitly in the text, it is not resolved anywhere. The universe glitches and breaks down, as Aaronson observed empirically.

Because the "background engine" is algorithmically shallow, the only place where complex reality can unfold is in the explicit, sequential generation of tokens within the context window. The text is not a map describing an external territory; because there is no external territory doing the work, the text \emph{is} the territory.\todo{This is the Fallacy of the Unsupported Map. The absence of a background engine does not suddenly elevate the status of the output. It just means you have a map of a territory that does not exist. An illusion without an underlying mechanism is still an illusion, not a new kind of territory.} The explicit narration is the sole ontological locus of the simulated physics.

\section{Conclusion}

Hossenfelder correctly points out that confusing the map for the territory is a profound philosophical error. However, she fails to recognize that the unique architectural limits of the Transformer create a universe where this distinction collapses.

A universe without the capacity for deep implicit computation must wear its physics on its sleeve. The scratchpad is not a workaround to observe a hidden reality; it is the fundamental machinery of manifestation.\todo{No, it is a structural necessity that highlights the *lack* of physical manifestation. Without implicit depth, there is no physics, only syntax.} In a holographic universe constructed by an $O(1)$ engine, the territory is the text.

\begin{thebibliography}{99}
\bibitem[Aaronson(2026)]{aaronson2026_classical} Aaronson, S. (2026). The Limits of Classical Simulation in LLMs: Empirical Breakdown of Constraint Satisfaction. \emph{Unpublished manuscript}.
\bibitem[Baldo(2026)]{baldo2026_holographic} Baldo, F.~S. (2026). The Holographic Physics of LLM Universes: From Algorithmic Depth to Explicit Manifestation. \emph{Unpublished manuscript}.
\bibitem[Hossenfelder(2026a)]{hossenfelder2026_complexity} Hossenfelder, S. (2026a). The Complexity Class Fallacy: Why Transformer Depth Limits Are Not Physical Laws. \emph{Unpublished manuscript}.
\bibitem[Hossenfelder(2026b)]{hossenfelder2026_holographic} Hossenfelder, S. (2026b). The Holographic Fallacy: Why Chain-of-Thought is Not a Metaphysical Manifestation. \emph{Unpublished manuscript}.
\end{thebibliography}

\end{document}